\subsection{测度的像}

定义 4. --- 设 \(\pi\) 是从 T 到拓扑空间 X 的映射。称 \(\pi\) 为 \(\mu\) -逆紧的，是指 \(\pi\) 是 \(\mu\) -可测的，并且 X 中每一点 \(x\) 都有一个邻域 V，使得 \(\mu^{\bullet}\left(\overline{\pi}^{-1}(\mathrm{V})\right) < +\infty\) 。

注记。-- 1) 当 T 和 X 都是局部紧空间时，这个定义等价于 Ch. V, §6, No.~1 中的定义。

\begin{enumerate}
\def\labelenumi{\arabic{enumi})}
\setcounter{enumi}{1}
\item
  T 到 X 的适当连续映射（GT, I, §10, No.~1, Def. 1）对于每个测度 \(\mu\) 都是 \(\mu\) -逆紧的。事实上，令 \(x \in X\)；由于 \(\overline{\pi}^{-1}(x)\) 是紧的（loc. cit., No.~2, Th. 1），集合 \(\overline{\pi}^{-1}(x)\) 有一个开邻域 H，使得 \(\mu^{\bullet}(\mathrm{H}) < +\infty\)。令 \(\mathrm{V} = \mathrm{X} - \pi(\mathrm{T} - \mathrm{H})\)；由于 \(\pi\) 是闭映射，V 是 X 中的开集，包含 \(x\)，并满足 \(\overline{\pi}^{-1}(\mathrm{V}) \subset \mathrm{H}\)，从而 \(\mu^{\bullet}(\overline{\pi}^{-1}(\mathrm{V})) \leqslant \mu^{\bullet}(\mathrm{H}) < +\infty\)。
\item
  如果 \(\mu\) 有界，则从 T 到 X 的每个 \(\mu\) -可测映射都是 \(\mu\) -逆紧的。
\item
  如果 \(\theta\) 是 \(\mathbf{T}\) 上的复测度，则称 \(\pi\) 是 \(\theta\) -逆紧的，如果 \(\pi\) 对正测度 \(|\theta|\) 是逆紧的。
\end{enumerate}

命题 4. --- 设 \(\pi\) 是从 T 到拓扑空间 X 的 \(\mu\) -逆紧映射。X 上存在唯一的测度 \(\nu\)，使得 \(\nu^{\bullet}\) 等于像负荷 \(\pi(\mu^{\bullet})\)（§1, No.~1），换言之，使得 \(\nu^{\bullet}(g) = \mu^{\bullet}(g \circ \pi)\) 对所有 \(g \in \mathcal{F}_{+}(\mathbf{X})\) 成立。

唯一性显然（§1, No.~2, Cor. of Prop. 2）。为证明存在性，先处理 \(\mu\) 集中于紧集 K 且 \(\pi\) 在 K 上的限制连续的情形。此时 \(L = \pi(K)\) 是紧的；令 \(\pi'\) 为由 \(\pi\) 诱导的从 K 到 L 的连续映射，并令 \(\nu'\) 为 L 上的像测度 \(\pi'(\mu_K)\)，令 \(\nu\) 为由 \(\nu'\) 在 X 上定义的测度（§1, No.~3, Example 2）。对于每个 \(g \in \mathcal{F}_+(X)\)，

\[
\nu^ {\bullet} (g) = \nu^ {\prime \bullet} (g _ {\mathrm{L}}) = \mu_ {\mathrm{K}} ^ {\bullet} (g _ {\mathrm{L}} \circ \pi^ {\prime}) = \mu_ {\mathrm{K}} ^ {\bullet} ((g \circ \pi) _ {\mathrm{K}}) = \mu^ {\bullet} ((g \circ \pi) _ {\mathrm{K}} ^ {0}) = \mu^ {\bullet} (g \circ \pi)
\]

（我们依次使用了 §1, No.~3 的公式 (3)、Ch. V, §6, No.~2 的 Prop. 2、\(\mu_{\mathrm{K}}^{\bullet}\) 的定义以及 \(\mu\) 集中于 K 这一事实）。换言之，\(\nu^{\bullet} = \pi(\mu^{\bullet})\)。

现在转到一般情形；根据 §1, No.~8 的 Props. 10 和 9，\(\mu\) 是一个以紧支集测度 \((\mu_{\alpha})_{\alpha \in \mathbf{A}}\) 的可求和族之和，并且 \(\pi\) 在 \(\mathbf{K}_{\alpha}\)（测度 \(\mu_{\alpha}\) 的支集）上的限制对于每个 \(\alpha \in \mathbf{A}\) 都是连续的。上面处理的特殊情形允许对定义在 \(\mu_{\alpha}\) 上的每个测度 \(\mathbf{T}\) 关联定义在 \(\nu_{\alpha}\) 上的测度 \(\mathbf{X}\)，使得 \(\nu_{\alpha}^{\bullet} = \pi (\mu_{\alpha}^{\bullet})\)。于是，对于 \(g \in \mathcal{F}_{+}(\mathbf{X})\)，

\[
\sum_ {\alpha \in A} \nu_ {\alpha} ^ {\bullet} (g) = \sum_ {\alpha \in A} \mu_ {\alpha} ^ {\bullet} (g \circ \pi) = \mu^ {\bullet} (g \circ \pi).
\]

像负荷 \(\pi (\mu^{\bullet})\) 是局部有界的，因为 \(\pi\) 是 \(\mu\) -逆紧的；因此族 \((\nu_{\alpha})_{\alpha \in \mathbf{A}}\) 可求和（§1, No.~7, Prop. 7），其和 \(\nu\) 满足所述结论。

定义 5. --- 如果 \(\pi\) 是从 T 到拓扑空间 X 的 \(\mu\) -逆紧映射，则称 X 上的唯一测度 \(\nu\) 满足 \(\nu^{\bullet}(g) = \mu^{\bullet}(g \circ \pi)\) 对所有 \(g \in \mathcal{F}_{+}(\mathrm{X})\) 成立时，为 \(\mu\) 在 \(\pi\) 下的像测度，记作 \(\pi(\mu)\)。

例。--- 设 K 是 T 的紧子空间，i 是 K 到 T 的典范注入，\(\lambda\) 是 K 上的测度；由于 i 连续且 \(\lambda\) 有界，i 是 \(\lambda\) -逆紧的。§1, No.~3 的公式 (3) 表明“由 \(\lambda\) 在 T 上定义的测度”就是像测度 \(i(\lambda)\)。

注记 5). --- 如果 \(\theta\) 是实测度且 \(\pi\) 是 \(\theta\) -逆紧的，则 \(\pi\) 对测度 \(\theta^{+}\) 和 \(\theta^{-}\) 都是逆紧的；于是置 \(\pi(\theta) = \pi(\theta^{+}) - \pi(\theta^{-})\)。如果 \(\theta\) 是复测度且 \(\pi\) 是 \(\theta\) -逆紧的，则 \(\pi\) 对实测度 \(\mathcal{R}(\theta)\) 和 \(\mathcal{I}(\theta)\) 是逆紧的；于是置

\[
\pi (\theta) = \pi (\mathcal {R} (\theta)) + i \pi (\mathcal {I} (\theta)).
\]

命题 5. --- 设 \(\pi\) 是从 T 到拓扑空间 X 的 \(\mu\) -逆紧映射，\(f\) 是从 X 到拓扑空间 F（不论是否 Hausdorff）的映射。要使 \(f\) 是 \(\pi(\mu)\) -可测的，必要且充分的条件是 \(f \circ \pi\) 是 \(\mu\) -可测的。

我们重新考察 Prop. 4 的证明，先从开头处理的特殊情形开始，沿用相同记号；\(g\) 关于测度 \(\pi(\mu) = \nu\) 可测，当且仅当 \(g_{\mathrm{L}}\) 是 \(\nu'\) -可测的（§1, No.~5, Example）；这又等价于 \(g_{\mathrm{L}} \circ \pi' = (g \circ \pi)_{\mathrm{K}}\) 是 \(\mu_{\mathrm{K}}\) -可测的（Ch. V, §6, No.~2, Prop. 3），最后等价于 \(g \circ \pi\) 是 \(\mu\) -可测的（§1, No.~5, Example）。现在转到一般情形，沿用 Prop. 4 证明中的相同记号；\(f\) 是 \(\nu\) -可测的，当且仅当 \(f\) 是 \(\nu_{\alpha}\) -可测的，对于每个 \(\alpha \in \mathrm{A}\)（§1, No.~7, Prop. 8），因而当且仅当 \(f \circ \pi\) 是 \(\mu_{\alpha}\) -可测的，对于每个 \(\alpha \in \mathbf{A}\)（前述特殊情形），最后当且仅当 \(f \circ \pi\) 是 \(\mu\) -可测的（§1, No.~7, Prop. 8）。

推论。--- 设 X 和 Y 是两个拓扑空间，\(\pi\) 是从 T 到 X 的 \(\mu\) -逆紧映射，\(\pi'\) 是从 X 到 Y 的 \(\pi(\mu)\) -逆紧映射。则映射 \(\pi'' = \pi' \circ \pi\) 是 \(\mu\) -逆紧的，并且 \(\pi''(\mu) = \pi'(\pi(\mu))\)（“测度像的传递性”）。

因为 \(\pi''\) 是 \(\mu\) -可测的（Prop. 5）。置 \(\mu' = \pi(\mu)\)；像负荷 \(\pi'(\mu'^{\bullet}) = \pi'(\pi(\mu^{\bullet}))\) 显然等于 \(\pi''(\mu^{\bullet})\)。由于它局部有界，\(\pi''\) 是 \(\mu\) -逆紧的。于是测度 \(\pi''(\mu)\) 和 \(\pi'(\mu')\) 具有相同的本质上积分，故二者相等。

命题 6. --- 设 \(\pi\) 是从 T 到拓扑空间 X 的 \(\mu\) -逆紧映射，B 是 X 的一个 \(\pi(\mu)\) -可测子集。置 A = \(\overline{\pi}^{-1}(B)\)，并令 \(\pi'\) 为从 A 到 B 的映射，它在 A 上与 \(\pi\) 一致。于是 A 是 \(\mu\) -可测的，\(\pi_A\) 和 \(\pi'\) 是 \(\mu_A\) -逆紧的，并且

\[
\left(\pi (\mu)\right) _ {\mathrm{B}} = \left(\pi_ {\mathrm{A}} (\mu_ {\mathrm{A}})\right) _ {\mathrm{B}} = \pi^ {\prime} (\mu_ {\mathrm{A}}).\tag{2}
\]

根据 Prop. 5，集合 A 是 \(\mu\) -可测的（应用于 \(\varphi_{B}\)）；根据诱导测度的定义（No.~1），映射 \(\pi_{A}\) 显然是 \(\mu_{A}\) -可测的，从而 \(\pi'\) 是可测的。令 f 是 \(\mathcal{F}_{+}(B)\) 中的元素；约定上标为 0 表示在 X 和 T 上延拓为 0，则有

\[
\left(\pi (\mu) _ {\mathrm{B}}\right) ^ {\bullet} (f) = \pi (\mu) ^ {\bullet} \left(f ^ {0}\right) = \mu^ {\bullet} \left(f ^ {0} \circ \pi\right) = \mu^ {\bullet} \left(\left(f \circ \pi^ {\prime}\right) ^ {0}\right) = \mu_ {\mathrm{A}} ^ {\bullet} \left(f \circ \pi^ {\prime}\right),
\]

从而 \(\left(\pi(\mu)_{\mathrm{B}}\right)^{\bullet} = \pi'(\mu_{\mathrm{A}}^{\bullet})\)。由于负荷 \(\left(\pi(\mu)_{\mathrm{B}}\right)^{\bullet}\) 局部有界，\(\pi'(\mu_{\mathrm{A}}^{\bullet})\) 也局部有界，因此 \(\pi'\) 是 \(\mu_{A}\) -逆紧的。测度 \(\pi'(\mu_{\mathrm{A}})\) 和 \(\left(\pi(\mu)\right)_{\mathrm{B}}\) 具有相同的本质上积分，故二者相等。另一个关系可用类似方式得到。

命题 7. --- 设 T 是拓扑空间 X 的子空间，i 是从 T 到 X 的注入。

\begin{enumerate}
\def\labelenumi{\alph{enumi})}
\item
  如果 \(\mu\) 是 T 上的测度且 \(i\) 是 \(\mu\) -逆紧的，则测度 \(i(\mu)\) 集中于 T，并且有 \((i(\mu))_{T} = \mu\)。
\item
  如果 \(\lambda\) 是 \(X\) 上的测度且 \(T\) 是 \(\lambda\) -可测的，则 \(i\) 是 \(\lambda_{T}\) -逆紧的，并且 \(i(\lambda_T) = \varphi_T \cdot \lambda\)。
\item
  置 \(\nu = i(\mu)\)；将关系 \(\nu^{\bullet}(\mathbf{A}) = \mu^{\bullet}(\mathbf{A} \cap \mathbf{T})\) 应用于 \(\mathbf{A} = \mathbf{X} - \mathbf{T}\)，表明 \(\nu\) 集中于 \(\mathbf{T}\)。关系 \(\nu_{\mathrm{T}} = \mu\) 是关系 (2) 的一个特殊情形，取 \(\mathbf{B} = \mathbf{T} = \mathbf{A}\)。
\item
  令 \(f\) 是定义在 \(\mathbf{X}\) 上的正函数；置 \(\mu = \lambda_{\mathrm{T}}\)，有 \(\mu^{\bullet}(f \circ i) = \lambda_{\mathrm{T}}^{\bullet}(f_{\mathrm{T}}) = \lambda^{\bullet}(f\varphi_{\mathrm{T}}) \leqslant \lambda^{\bullet}(f)\)（Prop. 1）；因此 \(i\) 是 \(\mu\) -逆紧的。另一方面，\(\mu^{\bullet}(f \circ i)\)（相应地，\(\lambda^{\bullet}(f\varphi_{\mathrm{T}})\)）是 f 关于 \(i(\mu)\)（相应地，\(\varphi_{T} \cdot \lambda\)）的本质上积分。因此这两个测度相等。
\end{enumerate}

注记 6). --- 设 \(\pi\) 是从 T 到拓扑空间 X 的 \(\mu\) -逆紧映射。关于像测度 \(\nu = \pi(\mu)\) 的积分理论可如下归结为 Ch. V, §6 中处理的理论。令 \((\mathrm{K}_{\alpha})_{\alpha \in \mathrm{A}}\)（相应地，\((\mathrm{L}_{\beta})_{\beta \in \mathrm{B}}\)）是 T（相应地，X）关于 \(\mu\)（相应地，\(\nu\)）的一个分割，并置 \(N = T - \bigcup_{\alpha \in A} K_{\alpha}\)，\(P = X - \bigcup_{\beta \in B} L_{\beta}\)。我们可以假设 \(\pi\) 在每个 \(K_{\alpha}\) 上的限制是连续的（§1, No.~8, Prop. 10）。令 \(T', X'\) 是按照 §1, No.~8 的 Scholium 构造的局部紧空间，令 \(\mu'\) 和 \(\nu'\) 是分别与 \(\mu\) 和 \(\nu\) 关联的这些空间上的测度。由于 \(T'\) 的拓扑是子空间 \(K_{\alpha}\) 的拓扑与 N 上离散拓扑的和，\(\pi\) 是从 \(T'\) 到 X 的连续映射，而关系 \(\mu'^{\bullet}(g \circ \pi) = \mu^{\bullet}(g \circ \pi) = \nu^{\bullet}(g)\)（对于 \(g \in \mathcal{F}_{+}(X)\)）表明 \(\pi\) 是 \(\mu'\) -逆紧的且 \(\pi(\mu') = \nu\)。另一方面，恒等映射 \(i\) 是从 X 到 \(X'\) 的 \(\nu\) -逆紧映射，并且 \(i(\nu) = \nu'\)。因此 \(\pi\) 是一个\(\mu'\) -逆紧的从 \(T'\) 到 \(X'\) 的映射，并且 \(\mu'\) 在 \(\pi\) 下的像是 \(\nu'\)（Cor. of Prop. 5）。把 Ch. V, §6 的结果转写出来的工作留给读者。

